\section{Results}
\subsection{Common support and spectral convergence}
The endpoint intersection contained 386,280 of 403,480 cells (95.74\%), excluding 17,200 cells from change denominators. Of the supported cells, 65,391 (16.93\%) met at least two adverse criteria and 5,493 (1.42\%) met all three (Fig.~\ref{fig:convergence}). The empirical thresholds were $\Delta$NDVI $\leq-0.00543$, $\Delta$MNDWI $\leq-0.03743$ and $\Delta$NDBI $\geq0.00835$. Convergence was therefore spatially selective. Its prevalence does not estimate the extent of damaged archaeology.

At 500 m, the NDVI- and NDBI-pressure ranks were strongly aligned ($\rho=0.972$), supporting their treatment within a shared surface-cover term. Five-epoch patterns varied across components and did not establish a uniform, monotonic deterioration trajectory. Complete epoch maps and threshold/harmonisation comparisons are supplied as supplementary evidence.

\begin{figure}[tbp]\centering
\includegraphics[width=.86\textwidth]{figure_04_spectral_pressure_convergence.pdf}
\caption{Convergence of adverse endpoint spectral criteria on the common E2004--E2024 mask. Values count coincident empirical criteria, with a denominator of 386,280 cells. The map is a landscape screen, not an observation of heritage damage.}\label{fig:convergence}
\end{figure}

\subsection{Priority depends on the uncertainty being considered}
Giri complex had the largest fixed 500 m score (0.591). For the separate integrated context score, the archived 50,000-draw weight analysis gave Giri a central 95\% rank interval of 1--9, with a top-three frequency of 0.774. Spatial resampling instead placed Bhallar at median rank 1 and Giri at median rank 2; the respective intervals were 1--5 and 1--3. These summaries answer different questions: the fixed score assumes the declared architecture, weight perturbation varies that architecture, and spatial resampling varies local spectral support while retaining terrain.

Rank concordance was higher between 500 and 1,000 m ($\rho=0.918$) than between 250 and 1,000 m ($\rho=0.711$). Removing whole domains produced substantial changes: landscape-only ranks correlated 0.325 with the full integrated reference, with a maximum movement of 11 positions; terrain-only ranks correlated 0.830, with a maximum movement of seven. Landscape-only scoring placed Sirkap first, whereas terrain-only scoring placed Bhallar first. A single leader therefore conceals materially different inspection rationales (Fig.~\ref{fig:ranks}).

\begin{figure}[tbp]\centering
\includegraphics[width=\textwidth]{figure_13_rank_uncertainty.pdf}
\caption{Archived component-rank sensitivity to weight and spatial perturbations. Intervals and top-set frequencies are conditional on their respective resampling designs. They do not combine every measurement, positional or modelling uncertainty. Full component tables and additional joint-uncertainty experiments accompany the notebook.}\label{fig:ranks}
\end{figure}

\subsection{Climate interpretation is product- and window-dependent}
Open-Meteo annual precipitation decreased by 10.30 mm year$^{-1}$ and annual mean temperature increased by 0.0418~$^\circ$C year$^{-1}$. The corrected residual-block intervals and annual-family adjusted probabilities are reported in Supplementary Table~S2. The point slopes were reproducible; their interpretation remained conditional on the selected gridded product.

NASA POWER gave opposing point slopes: +29.21 mm year$^{-1}$ for precipitation and $-0.0372$~$^\circ$C year$^{-1}$ for temperature. Annual precipitation agreement was weak across 35 years (Spearman $\rho=0.083$; RMSE 646.74 mm year$^{-1}$). High daily temperature rank agreement coexisted with a mean offset. These discrepancies reveal product sensitivity, without identifying either series as more accurate (Fig.~\ref{fig:weather}).

Acquisition matching also changed the epoch interpretation. E2024 mean antecedent precipitation was 42.47, 221.19 and 852.92 mm over 30, 90 and 365 days, respectively. It was not the wettest epoch under any of these windows: the maxima belonged to E2004, E2019 and E2014, respectively. A calendar-year or product-specific extreme therefore could not be assigned automatically to the satellite composite.

\begin{figure}[tbp]\centering
\includegraphics[width=\textwidth]{figure_07_weather_product_comparison.pdf}
\caption{Agreement between the two gridded weather products at matched temporal supports. Differences quantify product dependence, not error relative to station observations. Temporal correlation and absolute agreement are distinct.}\label{fig:weather}
\end{figure}

\subsection{Buffered model selection and diagnostics}
The stricter inner spatial comparison selected \ExtensionModel{} with mean development macro-F1 \ExtensionCV{}. Its outer macro-F1 was \ExtensionFOne{} (95\% test-block interval \ExtensionLow{}--\ExtensionHigh{}), and overall agreement was \ExtensionAgreement{}. Table~\ref{tab:models} reports all seven families, including those that did not lead either comparison. Paired differences, calibration and controls are provided in the Supplementary Information; the outer score was not used to replace the development-selected family.

\begin{table}[tbp]\centering\small
\caption{Seven-family spatial comparison. The primary model is selected using buffered development folds; outer intervals use 2,000 shared test-block resamples.}\label{tab:models}
\begin{tabular}{p{.35\linewidth}rrrr}\toprule
Model & Inner F1 & Outer F1 & 95\% interval & ECE \\\midrule
Random forest (primary) & 0.695 & 0.821 & 0.691--0.835 & 0.014 \\
Logistic regression & 0.655 & 0.807 & 0.680--0.822 & 0.050 \\
MLP & 0.624 & 0.849 & 0.723--0.888 & 0.020 \\
Extra trees & 0.617 & 0.826 & 0.694--0.840 & 0.029 \\
XGBoost & 0.616 & 0.827 & 0.699--0.845 & 0.017 \\
CatBoost & 0.614 & 0.792 & 0.680--0.825 & 0.017 \\
Histogram boosting & 0.611 & 0.846 & 0.704--0.857 & 0.051 \\
\bottomrule\end{tabular}
\end{table}


\ExtensionInterpretation{}

The selected random forest produced the same macro-F1 to numerical precision in Colab and Windows. XGBoost differed by 0.00061 between platforms (0.82687 versus 0.82748); this difference did not alter primary selection. The complete cross-platform comparison is archived.

The historical five-family reproduction yielded macro-F1 \LegacyFOne{}, agreement \LegacyAgreement{} and calibration error \LegacyECE{} for its development-selected MLP. These differ from the original stored targets of 0.830795, 0.833095 and 0.014280. The differences are preserved as historical reproduction discrepancies; targets were not reset to make a rerun pass. The stricter extension uses a different inner protocol and is reported as a separate experiment rather than a replacement claim of exact reproduction.
